\begin{mistake}{2026-04-22}{05}

\begin{problemcontent}
已知 3 阶矩阵 \(A\) 有三个特征值 \(-1, 2, \frac{1}{3}\)，对应的特征向量分别是 \(\xi_1, \xi_2, \xi_3\)，取 \(P=[2\xi_2, -\xi_1, 3\xi_3]\)，则 \(P^{-1}AP =\) \underline{\hspace{2cm}}。
\end{problemcontent}

\begin{solutioncontent}
设 \(P=[\eta_1, \eta_2, \eta_3]\)，其中 \(\eta_1=2\xi_2\)，\(\eta_2=-\xi_1\)，\(\eta_3=3\xi_3\)。
根据特征值定义 \(A\xi = \lambda\xi\)，特征向量数乘非零常数后，其对应的特征值不变：
\[
\begin{aligned}
A\eta_1 &= A(2\xi_2) = 2(A\xi_2) = 2(2\xi_2) = 2\eta_1 \\
A\eta_2 &= A(-\xi_1) = -(A\xi_1) = -(-1\xi_1) = -1\eta_2 \\
A\eta_3 &= A(3\xi_3) = 3(A\xi_3) = 3(\frac{1}{3}\xi_3) = \frac{1}{3}\eta_3
\end{aligned}
\]
利用分块矩阵乘法法则写成矩阵形式：
\[
AP = A[\eta_1, \eta_2, \eta_3] = [A\eta_1, A\eta_2, A\eta_3] = [2\eta_1, -1\eta_2, \frac{1}{3}\eta_3]
\]
提取右侧系数构造对角阵：
\[
AP = [\eta_1, \eta_2, \eta_3] \begin{bmatrix} 2 & 0 & 0 \\ 0 & -1 & 0 \\ 0 & 0 & \frac{1}{3} \end{bmatrix} = P \begin{bmatrix} 2 & 0 & 0 \\ 0 & -1 & 0 \\ 0 & 0 & \frac{1}{3} \end{bmatrix}
\]
两边同时左乘 \(P^{-1}\)，得：
\[
P^{-1}AP = \begin{bmatrix} 2 & 0 & 0 \\ 0 & -1 & 0 \\ 0 & 0 & \frac{1}{3} \end{bmatrix}
\]
\end{solutioncontent}

\mistakeknowledge{
\begin{itemize}
 \item \textbf{核心公式/定理}：若 \(A\xi = \lambda\xi\)，则对于常数 \(k \neq 0\)，\(A(k\xi) = \lambda(k\xi)\)，即倍乘不变 \(\lambda\)。
 \item \textbf{易错点}：误将特征向量的放大倍数乘到对角阵的特征值上；忽略了 \(P\) 矩阵列向量打乱了原特征值的排序。
 \item \textbf{关联知识}：相似对角化 \(P^{-1}AP=\Lambda\) 中，对角阵 \(\Lambda\) 元素的排列顺序必须与 \(P\) 中特征向量列的顺序严格一一对应。
\end{itemize}}

\end{mistake}
